\VerdierPhysicalPageAnchor{0021}
le morphisme de complexes de faisceaux :

\VerdierSemanticAnchor{intro-1-flat-quasi-isomorphism}
\[
F\otimes\rho:
F\otimes_{\mathbb Z}L_1^{\bullet}
\longrightarrow
F\otimes_{\mathbb Z}L_2^{\bullet}
\qquad\text{,}
\]
est un \emph{quasi-isomorphisme}.

\VerdierSemanticAnchor{intro-1-coherent-tensor-system}
On obtient donc un système cohérent de complexes de faisceaux
$F\otimes_{\mathbb Z}L^{\bullet}$ et de quasi-isomorphismes qui dépend, en un
sens qu'on peut préciser, fonctoriellement du faisceau $G$. En tout cas, les
objets de cohomologie des complexes $F\otimes_{\mathbb Z}L^{\bullet}$
($L^{\bullet}$ résolution plate de $G$) sont définis à isomorphismes uniques
près et dépendent fonctoriellement de $G$. Les foncteurs ainsi obtenus
possèdent les propriétés voulues et méritent le nom de foncteurs dérivés
gauches du foncteur ``produit tensoriel par $F$''.

\VerdierSemanticAnchor{intro-1-tor-sheaves}
Bien entendu, pour calculer les faisceaux de torsion
$\operatorname{Tor}^{\mathbb Z}(F,G)$, on peut prendre le faisceau associé au
préfaisceau
$U\longmapsto\operatorname{Tor}^{\mathbb Z}(F(U),G(U))$ ($U$ ouvert de $E$).
Mais la méthode que nous employons fournit un complexe de faisceaux bien défini
à quasi-isomorphisme près dont les objets de cohomologie sont
$\operatorname{Tor}^{\mathbb Z}(F,G)$ et $F\otimes_{\mathbb Z}G$, information
supplémentaire essentielle pour le maniement commode de ces foncteurs, ainsi
que nous le verrons par la suite. Cette méthode s'étend immédiatement au cas
des topos annelés et permet de définir dans ce cadre les foncteurs dérivés
gauches du foncteur ``produit tensoriel de faisceaux''.

\VerdierSemanticAnchor{intro-2}
\noindent\textbf{2.}\quad
Dans l'étude des foncteurs dérivés d'un foncteur composé de deux foncteurs, des
propriétés d'associativité, des relations du type de Künneth, on est amené à
étendre le formalisme des foncteurs dérivés au cas où l'argument n'est plus
seulement un objet de la catégorie étudiée, mais un complexe d'objets de cette
catégorie. Cette extension a été faite dans [1]. Rappelons-en les grandes
lignes. Soit $F:\mathcal A\to\mathcal B$ un foncteur additif entre deux
catégories abéliennes telles que $\mathcal A$ possède suffisamment d'objets
injectifs. Soit $Y^{\bullet}$ un complexe d'objets de $\mathcal A$. On définit
dans \emph{loc. cit.} les ``résolutions injectives'' du complexe $Y^{\bullet}$,
que nous appelons \emph{résolutions injectives de Cartan-Eilenberg} du complexe
$Y^{\bullet}$.

\VerdierSemanticAnchor{intro-2-ce-resolution}
Une telle résolution est, par définition, un double complexe
$I^{\bullet\bullet}$ dont les composants sont des objets injectifs de
$\mathcal A$, muni d'une augmentation :

\VerdierSemanticAnchor{intro-2-ce-augmentation}
\[
Y^{\bullet}\xrightarrow{\varepsilon}I^{\bullet\bullet}
\]
qui induit des résolutions injectives (au sens usuel) des composants, des
bords, et des objets de cohomologie du complexe $Y^{\bullet}$. Ces résolutions
sont uniques à homotopie (de doubles complexes) près. On forme alors le double
complexe $F(I^{\bullet\bullet})$, puis le complexe simple associé
$\int F(I^{\bullet\bullet})$ qui, lui aussi, est déterminé
